\section{\xdpninja: Non-Invasive XDP Capture}
\label{sec:xdp-ninja}

\xdpninja is the XDP host adapter for the \kunai library. It attaches
to running XDP programs at function-entry and function-exit boundaries
via BPF trampolines~\cite{hoiland2018xdp}, observes packets including
those the production program drops, and returns the production program
to its original behaviour on detach. The defining design constraint
is that the production XDP program is never modified, recompiled, or
restarted.

\subsection{Design Principle: Non-Invasive Attachment}
\label{sec:non-invasive}

Production XDP deployments at scale --- Cilium~\cite{cilium_docs},
Cloudflare's L4 load balancer~\cite{cloudflare_unimog}, and
Katran~\cite{shukla2020katran} --- represent significant operational
investment. Modifying or rebuilding them carries non-trivial risk:
data-path regression, downtime during reload, and re-validation across
the deployment fleet. The previously dominant XDP-native capture tool,
\texttt{xdpcap}~\cite{cloudflare2019xdpcap}, requires source-level
integration of capture hooks; this couples observability work to the
production deployment cycle and is rejected outright in operational
contexts where the production XDP source is not available to the
operator at all.

\xdpninja eliminates this friction with three properties. First,
production programs remain bit-identical between debugging sessions:
attach and detach do not touch the production binary. Second,
performance characteristics revert exactly when \xdpninja detaches,
because the trampoline disappears. Third, multiple \xdpninja
instances coexist on the same XDP program; runtime attach/detach is
the supported workflow.

The constraint shapes the rest of the architecture. \xdpninja cannot
call \texttt{bpf\_redirect\_map} (only valid from XDP itself), which
in turn rules out AF\_XDP zero-copy and cpumap dispatch as transports;
we discuss the resulting trade-off in \S\ref{sec:transport-tradeoff}.

\subsection{Three Attachment Modes}
\label{sec:attach-modes}

\xdpninja exposes three modes selected at the CLI:
\begin{description}
\item[\texttt{entry} (fentry):] observe a packet immediately before
the production XDP program runs.
\item[\texttt{exit} (fexit):] observe with the production program's
verdict (PASS, DROP, TX, REDIRECT) visible as an action atom in
\kunai's filter (e.g.\ \texttt{where action == XDP\_DROP}).
\item[\texttt{subfunction} (\texttt{--func}):] attach at any
\texttt{\_\_noinline} subfunction in the production XDP, observing
its intermediate state. The function is selected by name, resolved
at runtime through BTF.
\end{description}
The same kunai filter expression compiles for all three modes; only
the host adapter changes (\S\ref{sec:codegen}).

\subsection{BTF-Based Introspection}
\label{sec:btf-introspection}

Production XDP binaries embed BPF Type Format (BTF) information for
function signatures and subfunction layout. \xdpninja resolves
\texttt{--func} arguments and the trampoline attach target from the
running kernel's BTF representation of the loaded program; no
additional build step or sidecar metadata file is required.

\subsection{tcpdump-Compatible CLI}
\label{sec:cli}

The default output is pcapng to stdout, so \xdpninja drops directly
into existing capture pipelines:
\begin{lstlisting}[style=kunai,basicstyle=\footnotesize\ttfamily]
xdp-ninja -i eth0 \
   --dsl 'eth/ipv4/tcp where tcp.dport == 443' \
   | tshark -r -
\end{lstlisting}
\texttt{-w}~\textit{file} writes to disk; the binary accepts both kunai
DSL and (legacy) pcap-filter expressions, with kunai as the default.

\subsection{Beyond XDP: tc Mode}
\label{sec:tc-mode}

\kunai's compilation target is parametrised by a \texttt{Capabilities}
interface that abstracts host-specific knowledge: the action-fetcher
(XDP\_* vs TC\_ACT\_*) and the program-type prologue. We ship a tc
host adapter alongside the XDP one; it emits \texttt{SCHED\_CLS}
programs and attaches via TCX (Linux 6.6+) or via fentry/fexit on
existing tc BPF programs. Both attachment paths preserve the
non-invasive principle of \S\ref{sec:non-invasive}: nothing in the
production tc program is changed. We evaluate the tc host alongside
the XDP host in \S\ref{sec:eval-verifier}.

\TODO{Figure 3 (suggested): three attachment modes diagram. The
production XDP program is a single box; entry / subfunction / exit
hooks are arrows that intercept it without modifying the box. Useful
to anchor \S\ref{sec:attach-modes} visually.}
